% ============================================================
\setpartpreamble{\partintro
  {Metadaten, Fluent API, Conventions, Validierung und neue Spaltentypen.}
  {
    \item Kapitel~\ref{ch:annotations-and-fluent-api}: Annotations und Fluent API
    \item Kapitel~\ref{ch:conventions}: Conventions
    \item Kapitel~\ref{ch:model-validation}: Modellvalidierung
    \item Kapitel~\ref{ch:type-mappings-and-value-converters}: Type Mappings und Value
    Converter
  }}
\part{Das Modell}\label{part:III}

% ------------------------------------------------------------
\chapter{Annotations und Fluent API}\label{ch:annotations-and-fluent-api}

Annotations sind EFs Mechanismus für erweiterbare Metadaten: Schlüssel-Wert-Paare an jedem
Modellelement (Modell, Entity-Typ, Property, Index, ...). Ein Plugin speichert seine
Konfiguration als Annotations, und die gesamte restliche Infrastruktur (Conventions,
Migrations, Snapshots, Compiled Models) kann damit arbeiten.

\section{Regeln für Annotations}

\begin{datatable}{H{0.75}H{1.25}}
  \textbf{Regel} & \textbf{Begründung} \\
  \tablehead
  Eigenes Präfix (\code{Acme:Rls:}) & Vermeidet Kollisionen mit EF, Providern und anderen
  Plugins. \\
  \tablerule
  Nur einfache Werte: string, bool, Zahlen, Enums & Snapshots und Migrationen schreiben
  Werte als C\#-Literale. Eigene Klassen lassen sich nicht serialisieren. \\
  \tablerule
  Komplexe Werte als String (z. B. JSON) & Strings vergleicht der Differ zuverlässig per
  \code{Equals}. Ein Record mit \code{string[]} würde per Referenz verglichen und ständig
  „geändert“ wirken. \\
  \tablerule
  Default-Werte nicht speichern & „Annotation fehlt“ und „Annotation hat Default“ sind für
  den Differ verschieden und erzeugen leere Migrationen. \\
  \tablerule
  Je Teilkonfiguration eine Annotation & Jede Policy hat eine eigene Annotation. Dann zeigt
  der Diff genau, welche sich geändert hat. \\
\end{datatable}

\begin{sourcecode}[Metadata/\allowbreak{}Rls\allowbreak{}Annotation\allowbreak{}Names.\allowbreak{}cs,\allowbreak{} Metadata/\allowbreak{}Rls\allowbreak{}Policy.\allowbreak{}cs (\allowbreak{}gekürzt)]{csharp}
public static class RlsAnnotationNames
{
    public const string Prefix = "Acme:Rls:";
    public const string Enabled = Prefix + "Enabled";
    public const string Forced = Prefix + "Forced";
    public const string PolicyPrefix = Prefix + "Policy:";     // + Policy-Name
    public const string QueryFilterKey = "Acme.Tenancy";
}

public sealed record RlsPolicy(string Name, PolicyCommand Command, bool Permissive,
    string[]? Roles, string? Using, string? WithCheck)
{
    public string Serialize()
        => JsonSerializer.Serialize(this, RlsJsonContext.Default.RlsPolicy);

    public static RlsPolicy Deserialize(string json) => /* ... */;
}
\end{sourcecode}

\section{Die Fluent API}

Fluent-API-Methoden sind Extension-Methods auf den Buildern von EF. Drei Gestaltungsregeln
machen sie professionell: Die \textbf{nicht-generische} Variante ist die eigentliche
Implementierung (der Model-Snapshot nutzt nicht-generische Builder,
Kapitel~\ref{ch:design-time-services}). Die \textbf{generische} Variante delegiert nur und
erhält den Typ für weiteres Chaining. Und die Klasse liegt im Namespace
\code{Microsoft.\allowbreak{}Entity\allowbreak{}Framework\allowbreak{}Core}, damit sie
überall ohne \code{using} verfügbar ist.

\begin{sourcecode}[Builders/\allowbreak{}Rls\allowbreak{}Entity\allowbreak{}Type\allowbreak{}Builder\allowbreak{}Extensions.\allowbreak{}cs (\allowbreak{}gekürzt)]{csharp}
namespace Microsoft.EntityFrameworkCore;

public static class RlsEntityTypeBuilderExtensions
{
    public static EntityTypeBuilder HasRowLevelSecurity(
        this EntityTypeBuilder builder, bool force = false)
    {
        ArgumentNullException.ThrowIfNull(builder);
        builder.Metadata.SetAnnotation(RlsAnnotationNames.Enabled, true);

        if (force)
            builder.Metadata.SetAnnotation(RlsAnnotationNames.Forced, true);
        else
            builder.Metadata.RemoveAnnotation(RlsAnnotationNames.Forced);   // kein Default!

        return builder;
    }

    public static EntityTypeBuilder HasPolicy(
        this EntityTypeBuilder builder, string name, Action<PolicyBuilder> configure)
    {
        var policyBuilder = new PolicyBuilder(name);
        configure(policyBuilder);
        builder.Metadata.SetAnnotation(
            RlsAnnotationNames.PolicyPrefix + name, policyBuilder.Build().Serialize());
        return builder;
    }

    public static EntityTypeBuilder<TEntity> HasRowLevelSecurity<TEntity>(
        this EntityTypeBuilder<TEntity> builder, bool force = false) where TEntity : class
        => (EntityTypeBuilder<TEntity>)HasRowLevelSecurity(
            (EntityTypeBuilder)builder, force);

    // ... HasNoRowLevelSecurity, HasPolicy<TEntity> analog
}
\end{sourcecode}

Verschachtelte Konfiguration wie bei \code{Has\allowbreak{}Policy} nimmt ein
\code{Action<\allowbreak{}Policy\allowbreak{}Builder>} entgegen. Der Builder prüft die
Regeln der Datenbank schon beim Modellbau. Ein Fehler wie „INSERT-Policy mit USING“ fällt so
beim ersten Start der Anwendung auf und nicht erst beim Ausführen der Migration in
Produktion:

\begin{sourcecode}[Builders/\allowbreak{}Policy\allowbreak{}Builder.\allowbreak{}cs › Build]{csharp}
internal RlsPolicy Build()
{
    if (_using is null && _withCheck is null)
        throw new InvalidOperationException(
            $"Policy '{_name}' braucht USING oder WITH CHECK.");

    if (_command == PolicyCommand.Insert && _using is not null)
        throw new InvalidOperationException(
            $"INSERT-Policy '{_name}' erlaubt nur WITH CHECK.");

    if (_command is PolicyCommand.Select or PolicyCommand.Delete && _withCheck is not null)
        throw new InvalidOperationException(
            $"{_command}-Policy '{_name}' erlaubt nur USING.");

    return new RlsPolicy(_name, _command, _permissive,
        _roles.Count > 0 ? [.. _roles] : null, _using, _withCheck);
}
\end{sourcecode}

\section{Konfigurationsquellen}

EF merkt sich zu jedem Metadatum, \emph{woher} es stammt. Eine Quelle kann eine gleich
starke oder schwächere überschreiben, aber nie eine stärkere. Das ist der Mechanismus, mit
dem Plugins sinnvolle Defaults setzen, die Nutzer jederzeit übersteuern können:

\begin{datatable}{H{0.82}H{1.64}H{0.55}}
  \textbf{ConfigurationSource} & \textbf{Gesetzt durch} & \textbf{Stärke} \\
  \tablehead
  \code{Explicit} & Fluent API im \code{On\allowbreak{}Model\allowbreak{}Creating}
  (\code{Metadata.\allowbreak{}Set\allowbreak{}Annotation}) & stärkste \\
  \tablerule
  \code{Data\allowbreak{}Annotation} & Attribut-Conventions
  (\code{from\allowbreak{}Data\allowbreak{}Annotation: true}) & mittel \\
  \tablerule
  \code{Convention} & Conventions über \code{IConvention*Builder} & schwächste \\
\end{datatable}

Für eigene Conventions bietet das Plugin deshalb zusätzlich Überladungen auf dem
Convention-Builder an. Sie geben \code{null} zurück, wenn eine stärkere Quelle anders
entschieden hat:

\begin{sourcecode}[Builders/\allowbreak{}Rls\allowbreak{}Entity\allowbreak{}Type\allowbreak{}Builder\allowbreak{}Extensions.\allowbreak{}cs › Convention-Variante]{csharp}
public static IConventionEntityTypeBuilder? HasRowLevelSecurity(
    this IConventionEntityTypeBuilder builder,
    bool enabled,
    bool fromDataAnnotation = false)
    => builder.HasAnnotation(RlsAnnotationNames.Enabled, enabled, fromDataAnnotation)
           is null
        ? null
        : builder;
\end{sourcecode}

Die Prüfung der Quelle steckt im Builder, nicht in den Metadaten. Wer in einer Convention
direkt an die Metadaten schreibt, umgeht sie:

\begin{sourcecode}[Convention (\allowbreak{}FALSCH)]{csharp}
// In ProcessModelFinalizing: direkt an den Metadaten statt über entityType.Builder
entityType.SetAnnotation(RlsAnnotationNames.Enabled, true);
\end{sourcecode}

\begin{failure}[Was passiert, wenn eine Convention die Konfigurationsquelle umgeht?]
  \textbf{Beobachtbar.} Der Anwender schaltet RLS für eine Entität mit
  \code{Has\allowbreak{}No\allowbreak{}Row\allowbreak{}Level\allowbreak{}Security(\allowbreak{})}
  ab, und nach dem Modellbau ist es wieder an. Die nächste Migration aktiviert RLS für eine
  Tabelle, die ausdrücklich ausgenommen war.

  \textbf{Ursache.} \code{Set\allowbreak{}Annotation} am Convention-Metadatum überschreibt
  jeden vorhandenen Wert. Nur der Convention-Builder vergleicht die Quellen und gibt
  \code{null} zurück, wenn eine stärkere entschieden hat.

  \textbf{Lösung.} In Conventions immer über \code{entity\allowbreak{}Type.\allowbreak{}Builder}
  schreiben oder die Convention-Überladungen des Plugins verwenden.

  \textbf{Test:}
  \code{Failure\allowbreak{}Case\allowbreak{}Tests.\allowbreak{}Convention\allowbreak{}\_writing\allowbreak{}\_to\allowbreak{}\_the\allowbreak{}\_metadata\allowbreak{}\_…}
  und die Gegenprobe über den Builder.
\end{failure}

% ------------------------------------------------------------
\chapter{Conventions}\label{ch:conventions}

Conventions verändern das Modell automatisch. EF selbst besteht zum größten Teil aus
Conventions: Schlüssel erkennen, Beziehungen finden, Tabellennamen ableiten. Ein Plugin fügt
eigene hinzu, um Konfiguration überflüssig zu machen.

\section{Convention-Typen und Zeitpunkte}

Eine Convention implementiert ein oder mehrere Interfaces. Jedes steht für ein Ereignis
während des Modellbaus. Die wichtigsten:

\begin{datatable}{H{1.04}H{0.91}H{1.04}}
  \textbf{Interface} & \textbf{Ereignis} & \textbf{Typische Verwendung} \\
  \tablehead
  \code{IEntity\allowbreak{}Type\allowbreak{}Added\allowbreak{}Convention} & Entity-Typ
  wurde hinzugefügt & Attribute auswerten, Grundkonfiguration \\
  \tablerule
  \code{IProperty\allowbreak{}Added\allowbreak{}Convention} & Property wurde hinzugefügt &
  Spaltentypen, Namenskonventionen \\
  \tablerule
  \code{IForeign\allowbreak{}Key\allowbreak{}Added\allowbreak{}Convention} & Beziehung
  entstanden & Indizes, Löschverhalten \\
  \tablerule
  \code{IEntity\allowbreak{}Type\allowbreak{}Base\allowbreak{}Type\allowbreak{}Changed\allowbreak{}Convention}
  & Vererbung geändert & Konfiguration verschieben \\
  \tablerule
  \code{IModel\allowbreak{}Finalizing\allowbreak{}Convention} & Modell wird abgeschlossen
  (noch änderbar) & Alles, was vom fertigen Modell abhängt \\
  \tablerule
  \code{IModel\allowbreak{}Finalized\allowbreak{}Convention} & Modell ist fertig (read-only)
  & Validierung \\
\end{datatable}

Der häufigste Convention-Fehler ist der falsche Zeitpunkt:

\begin{sourcecode}[Convention (\allowbreak{}FALSCH)]{csharp}
public void ProcessPropertyAdded(IConventionPropertyBuilder propertyBuilder,
    IConventionContext<IConventionPropertyBuilder> context)
{
    // Spaltennamen gleich merken, etwa für den Policy-Ausdruck
    propertyBuilder.HasAnnotation("Acme:Column", propertyBuilder.Metadata.GetColumnName());
}
\end{sourcecode}

\begin{failure}[Was passiert, wenn eine Convention zu früh läuft?]
  \textbf{Beobachtbar.} Konfiguriert der Anwender danach
  \code{Has\allowbreak{}Column\allowbreak{}Name(\textquotedbl{}tenant\textquotedbl{})} oder
  benennt ein Naming-Plugin die Spalte um, bleibt der gemerkte Name
  \code{Tenant\allowbreak{}Id}. Die Policy verweist auf eine Spalte, die es nicht gibt, und
  die Migration scheitert erst in der Datenbank.

  \textbf{Ursache.} Reaktive Conventions (\code{*Added}, \code{*Changed}) laufen, während
  der Anwender noch konfiguriert. Ein Spaltenname kann sich danach noch ändern, ein
  Entity-Typ noch einen Basistyp bekommen.

  \textbf{Lösung.} Was vom fertigen Modell abhängt, in einer
  \code{IModel\allowbreak{}Finalizing\allowbreak{}Convention} lesen, und zwar über
  \code{Get\allowbreak{}Column\allowbreak{}Name(\allowbreak{}table)}. Wer reaktiv bleiben
  muss, reagiert auf \emph{alle} relevanten Ereignisse.

  \textbf{Test:}
  \code{Failure\allowbreak{}Case\allowbreak{}Tests.\allowbreak{}Convention\allowbreak{}\_reading\allowbreak{}\_the\allowbreak{}\_column\allowbreak{}\_name\allowbreak{}\_too\allowbreak{}\_early\allowbreak{}\_…}
  und die Gegenprobe bei \code{Model\allowbreak{}Finalizing}.
\end{failure}

\section{Conventions einhängen}

Das Plugin-Interface \code{IConvention\allowbreak{}Set\allowbreak{}Plugin} bekommt das
Convention-Set und fügt eigene Conventions hinzu.
\code{Convention\allowbreak{}Set.\allowbreak{}Add} trägt eine Convention in alle Listen ein,
deren Interfaces sie implementiert, jeweils am Ende:

\begin{sourcecode}[Conventions/\allowbreak{}Tenancy\allowbreak{}Convention\allowbreak{}Set\allowbreak{}Plugin.\allowbreak{}cs]{csharp}
public class TenancyConventionSetPlugin : IConventionSetPlugin
{
    private readonly ProviderConventionSetBuilderDependencies _dependencies;
    private readonly IDbContextOptions _options;
    private readonly ISqlGenerationHelper _sqlGenerationHelper;
    private readonly TenancyLoggingDefinitions _loggingDefinitions;

    public TenancyConventionSetPlugin(
        ProviderConventionSetBuilderDependencies dependencies,
        IDbContextOptions options,
        ISqlGenerationHelper sqlGenerationHelper,
        TenancyLoggingDefinitions loggingDefinitions)
    {
        // ... ArgumentNullException.ThrowIfNull für alle Parameter ...
        _dependencies = dependencies;
        _options = options;
        _sqlGenerationHelper = sqlGenerationHelper;
        _loggingDefinitions = loggingDefinitions;
    }

    public virtual ConventionSet ModifyConventions(ConventionSet conventionSet)
    {
        var settings = _options.FindExtension<TenancyOptionsExtension>()?.Settings;
        if (settings is null)
            return conventionSet;

        // Kapitel „Value Generation“
        conventionSet.Add(new TenantIdKeyValueGenerationConvention());
        conventionSet.Add(new NoTenantIsolationAttributeConvention(_dependencies));
        conventionSet.Add(new TenantIsolationConvention(
            settings, _sqlGenerationHelper, _dependencies.Logger, _loggingDefinitions));
        conventionSet.Add(new TenancyModelValidationConvention(
            _dependencies.Logger, _loggingDefinitions));

        return conventionSet;
    }
}
\end{sourcecode}

\code{Provider\allowbreak{}Convention\allowbreak{}Set\allowbreak{}Builder\allowbreak{}Dependencies}
ist das Parameterobjekt, das auch EFs eigene Conventions bekommen. Es enthält unter anderem
den Model-Logger, den Validation-Logger und die Type-Mapping-Source. Andere Dienste wie
\code{ISql\allowbreak{}Generation\allowbreak{}Helper} lassen sich zusätzlich injizieren,
weil das Plugin selbst im internen Provider aufgelöst wird.

\section{Die Isolations-Convention}

Die zentrale Convention des Beispiels läuft bei \code{Model\allowbreak{}Finalizing}, weil
Spaltennamen erst dann feststehen. Sie findet die Mandanten-Property (per
\code{Is\allowbreak{}Tenant\allowbreak{}Property(\allowbreak{})} markiert oder
\code{Tenant\allowbreak{}Id} bei \code{ITenant\allowbreak{}Owned},
Kapitel~\ref{ch:property-level-fluent-api-and}), aktiviert RLS mit Policy, legt einen Index
an, setzt eine Laufzeit-Annotation für den SaveChanges-Interceptor und optional einen
Query-Filter:

\begin{sourcecode}[Conventions/\allowbreak{}Tenant\allowbreak{}Isolation\allowbreak{}Convention.\allowbreak{}cs › Process\allowbreak{}Model\allowbreak{}Finalizing]{csharp}
public virtual void ProcessModelFinalizing(
    IConventionModelBuilder modelBuilder,
    IConventionContext<IConventionModelBuilder> context)
{
    // TPT, TPC, Entity Splitting, Owned Types in eigener Tabelle und implizite
    // Join-Tabellen ließen Mandantendaten ohne Mandantenspalte zurück: zuerst ablehnen
    TenantTableMappingValidator.Validate(modelBuilder.Metadata, e => e.IsRlsOptedOut());

    foreach (var entityType in modelBuilder.Metadata.GetEntityTypes())
    {
        if (entityType.BaseType is not null       // TPH: Tabelle des Root-Typs
            || entityType.IsOwned())              // geprüft: Eigentümertabelle oder JSON
            continue;

        // Markierte Property oder "TenantId" bei ITenantOwned (Kernpaket,
        // Kapitel „Fluent API auf Property-Ebene und Laufzeit-Annotations“)
        if (entityType.FindTenantPropertyCandidate() is not IConventionProperty property)
            continue;                             // (+ Typprüfung, gekürzt)

        if (entityType.IsRlsOptedOut())
        {
            _logger.TenantIsolationOptedOut(_loggingDefinitions, entityType);  // Warnung
            continue;
        }

        // Keine Tabelle (View, ToSqlQuery, Keyless) -> kein RLS
        if (StoreObjectIdentifier.Create(entityType, StoreObjectType.Table)
            is not { } table)
            continue;

        // Laufzeit-Annotation: der SaveChanges-Interceptor findet so die Property
        entityType.Builder.HasAnnotation(
            TenancyAnnotationNames.TenantProperty, property.Name);

        ApplyRowLevelSecurity(entityType, property, table);
        if (_settings.CreateTenantIndex) EnsureTenantIndex(entityType, property);
        if (_settings.ApplyQueryFilter) ApplyQueryFilter(entityType, property);

        _logger.TenantIsolationApplied(_loggingDefinitions, entityType);       // Debug
    }
}
\end{sourcecode}

\begin{sourcecode}[Conventions/\allowbreak{}Tenant\allowbreak{}Isolation\allowbreak{}Convention.\allowbreak{}cs › Apply\allowbreak{}Row\allowbreak{}Level\allowbreak{}Security]{csharp}
protected virtual void ApplyRowLevelSecurity(
    IConventionEntityType entityType,
    IConventionProperty property,
    StoreObjectIdentifier table)
{
    // Spaltenname für genau diese Tabelle (bei Table Splitting kann er abweichen)
    var column = _sqlGenerationHelper.DelimitIdentifier(property.GetColumnName(table)!);
    var condition =
        $"{column} = nullif(current_setting('{_settings.SettingName}', true), '')::uuid";

    var policy = new RlsPolicy(_settings.PolicyName, PolicyCommand.All, Permissive: true,
        Roles: null, Using: condition, WithCheck: condition);

    // Convention-Builder: überschreibt explizite Konfiguration NICHT
    var builder = entityType.Builder;
    builder.HasAnnotation(RlsAnnotationNames.Enabled, true);
    if (_settings.ForceRowLevelSecurity)
        builder.HasAnnotation(RlsAnnotationNames.Forced, true);
    builder.HasAnnotation(RlsAnnotationNames.PolicyPrefix + _settings.PolicyName,
        policy.Serialize());
}
\end{sourcecode}

\begin{important}[Selbst erzeugtes SQL validieren]
  Der Setting-Name wird unmaskiert in den Policy-Ausdruck geschrieben. Anders als bei SQL,
  das der Entwickler selbst in
  \code{Using(\allowbreak{}.\allowbreak{}.\allowbreak{}.\allowbreak{})} schreibt, setzt hier
  das \emph{Plugin} SQL zusammen und ist dafür verantwortlich, dass nichts Ungültiges
  hineinkommt. Deshalb validiert der Options-Builder den Namen mit einem strengen regulären
  Ausdruck
  (\code{[a-z\_\allowbreak{}][a-z0-9\_\allowbreak{}]*\textbackslash{}.\allowbreak{}[a-z\_\allowbreak{}][a-z0-9\_\allowbreak{}]*}).
\end{important}

\section{Welche Tabellenabbildungen geschützt sind}\label{sec:supported-table-mappings}

RLS schützt Tabellen, keine Entitäten. Die Policy vergleicht die Mandantenspalte \emph{derselben}
Tabelle mit dem Session-Setting. Eine Tabelle mit Mandantendaten, aber ohne diese Spalte, lässt
sich mit dieser Policy nicht schützen. Eine Entität kann mehrere Tabellen belegen, und mehrere
Entitäten können sich eine Tabelle teilen. Deshalb prüft
\code{Tenant\allowbreak{}Table\allowbreak{}Mapping\allowbreak{}Validator} aus dem Kernpaket
vor jeder Annotation, ob jede Tabelle mit Mandantendaten die Mandantenspalte enthält. Beide
Provider-Pakete rufen ihn auf, der Vertrag ist also für PostgreSQL und SQL Server gleich:

\begin{datatable}{H{1.00}H{1.25}H{0.75}}
  \textbf{Abbildung} & \textbf{Tabellen mit Mandantendaten} & \textbf{Verhalten} \\
  \tablehead
  Eine Tabelle pro Entität & die Tabelle, mit Mandantenspalte & geschützt \\
  \tablerule
  TPH (Standard) & eine Tabelle für die ganze Hierarchie & geschützt, Policy am Root-Typ \\
  \tablerule
  Mandantenbezogener abgeleiteter Typ unter einem Root ohne Mandant & dieselbe Tabelle, Zeilen
  anderer Typen ohne Mandant & abgelehnt \\
  \tablerule
  TPT & je abgeleitetem Typ eine Tabelle \emph{ohne} Mandantenspalte & abgelehnt \\
  \tablerule
  TPC & je konkretem Typ eine Tabelle, die abstrakte Basis hat keine & abgelehnt (nicht
  implementiert) \\
  \tablerule
  Owned Type in der Eigentümertabelle (\code{Owns\allowbreak{}One}) & Spalten der
  Eigentümertabelle & geschützt \\
  \tablerule
  Owned Type als JSON (\code{To\allowbreak{}Json}) & eine Spalte der Eigentümertabelle &
  geschützt \\
  \tablerule
  Owned Type in eigener Tabelle (\code{Owns\allowbreak{}Many}, \code{To\allowbreak{}Table})
  & eigene Tabelle ohne Mandantenspalte & abgelehnt \\
  \tablerule
  Table Splitting & dieselben Zeilen für alle Typen & geschützt, wenn alle mandantenbezogenen
  Typen dieselbe Spalte nutzen \\
  \tablerule
  Entity Splitting (\code{Split\allowbreak{}To\allowbreak{}Table}) & weitere Tabelle ohne
  Mandantenspalte & abgelehnt \\
  \tablerule
  Implizite Many-to-many-Tabelle mit mandantenbezogener Seite & Join-Tabelle ohne
  Mandantenspalte & abgelehnt \\
  \tablerule
  Join-Entität mit Mandanten-Property (\code{Using\allowbreak{}Entity}) & Join-Tabelle mit
  Mandantenspalte & geschützt \\
  \tablerule
  View, \code{To\allowbreak{}Sql\allowbreak{}Query}, keyless ohne Tabelle & keine eigene
  Tabelle & keine Policy; die Basistabellen tragen sie \\
\end{datatable}

Abgelehnt heißt: Das Modell wirft beim ersten Zugriff auf \code{Model} eine
\code{Invalid\allowbreak{}Operation\allowbreak{}Exception}, die Tabelle, Typ und Ausweg nennt.
Das passiert vor \code{Ensure\allowbreak{}Created}, vor jeder Migration und vor der ersten
Abfrage. Eine scheinbare Unterstützung, die eine Policy auf eine Tabelle ohne passende Spalte
schreibt, wäre schlimmer als die Ablehnung. TPC wäre grundsätzlich schützbar, weil jede
konkrete Tabelle die Mandantenspalte erbt. Das Plugin unterstützt aber genau eine Tabelle pro
Hierarchie, und eine halbe Umsetzung soll keine Sicherheit vortäuschen. Als Ausweg bleiben TPH
oder eigenständige Entitäten.

Bewusste Ausnahmen bleiben möglich: Eine mit \code{[No\allowbreak{}Tenant\allowbreak{}Isolation]}
oder \code{Has\allowbreak{}No\allowbreak{}Row\allowbreak{}Level\allowbreak{}Security(\allowbreak{})}
ausgenommene Entität wird nicht geprüft, und die Convention protokolliert die Ausnahme als
Warnung. So unterscheidet das Modell eine Entscheidung von einem Versehen. Für die implizite
Join-Tabelle gibt es im PostgreSQL-Paket
\code{Using\allowbreak{}Entity(\allowbreak{}j => j.\allowbreak{}Has\allowbreak{}No\allowbreak{}Row\allowbreak{}Level\allowbreak{}Security(\allowbreak{}))};
das SQL-Server-Paket kennt nur das Attribut und braucht dafür eine Join-Klasse. Besser ist
fast immer eine Mandanten-Property an der Join-Entität:

\begin{sourcecode}[Anwendung]{csharp}
modelBuilder.Entity<Basket>()
    .HasMany(b => b.Labels).WithMany(l => l.Baskets)
    .UsingEntity("BasketLabel",
        j => j.Property<TenantId>("TenantId").IsTenantProperty());   // Shadow-Property
\end{sourcecode}

Die Shadow-Property füllt der SaveChanges-Interceptor wie jede andere Mandanten-Property.
\code{Mapping\allowbreak{}Isolation\allowbreak{}Tests} prüft beides gegen ein echtes
PostgreSQL mit eingeschränkter Rolle: Ein TPC-Modell scheitert, bevor eine einzige Tabelle
existiert. In einem Modell mit TPH, Owned Type, JSON-Spalte und dieser Join-Entität hat jede
Tabelle RLS, und ohne Mandant ist in keiner Tabelle eine Zeile sichtbar.
\code{Table\allowbreak{}Mapping\allowbreak{}Tests} und
\code{Sql\allowbreak{}Server\allowbreak{}Table\allowbreak{}Mapping\allowbreak{}Tests} decken
jede Zeile der Matrix ohne Datenbank ab.

\section{Strukturelle Änderungen: Index und Query-Filter}

Conventions dürfen nicht nur Annotations setzen, sondern das Modell strukturell verändern.
Der Index wird nur angelegt, wenn nicht schon ein Index oder Schlüssel mit der
Mandantenspalte beginnt. Der Query-Filter wird als Expression Tree gebaut und als
\textbf{benannter Filter} registriert, eine Neuerung in EF Core 10. Benannte Filter
koexistieren mit eigenen Filtern des Nutzers und lassen sich gezielt per
\code{Ignore\allowbreak{}Query\allowbreak{}Filters(\allowbreak{}[\textquotedbl{}Acme.\allowbreak{}Tenancy\textquotedbl{}])}
abschalten.

\begin{sourcecode}[Conventions/\allowbreak{}Tenant\allowbreak{}Isolation\allowbreak{}Convention.\allowbreak{}cs (\allowbreak{}Auszug)]{csharp}
protected virtual void EnsureTenantIndex(
    IConventionEntityType entityType, IConventionProperty property)
{
    var covered = entityType.GetIndexes().Any(i => i.Properties[0] == property)
                  || entityType.GetKeys().Any(k => k.Properties[0] == property);
    if (!covered)
        entityType.Builder.HasIndex(new[] { property });
}

protected virtual void ApplyQueryFilter(
    IConventionEntityType entityType, IConventionProperty property)
{
    // e => EF.Property<TenantId>(e, "...") == EF.Functions.CurrentTenant()
    // EF.Property funktioniert für jeden Property-Namen und auch für Shadow-Properties
    var parameter = Expression.Parameter(entityType.ClrType, "e");
    var body = Expression.Equal(
        Expression.Call(typeof(EF), nameof(EF.Property), [typeof(TenantId)],
            parameter, Expression.Constant(property.Name)),
        Expression.Call(CurrentTenantMethod, Expression.Property(null, FunctionsProperty)));

    entityType.Builder.HasQueryFilter(
        RlsAnnotationNames.QueryFilterKey, Expression.Lambda(body, parameter));
}
\end{sourcecode}

Bemerkenswert ist, dass der Filter die eigene Funktion
\code{EF.\allowbreak{}Functions.\allowbreak{}Current\allowbreak{}Tenant(\allowbreak{})}
verwendet. Er wird damit über den Translator aus
Kapitel~\ref{ch:translating-methods-and-members} nach SQL übersetzt. Die Bausteine des
Plugins greifen ineinander.

\section{Endgültige Namen, egal wer sie setzt}

Die Isolations-Convention liest den Spaltennamen erst bei
\code{Model\allowbreak{}Finalizing} und nur über
\code{Get\allowbreak{}Column\allowbreak{}Name(\allowbreak{}table)}. Namen, die andere
Conventions oder Plugins ändern, stehen zu diesem Zeitpunkt fest. Das Beispiel prüft das
gegen EFCore.NamingConventions, das Tabellen und Spalten in snake\_case umbenennt, und zwar
in beiden Aktivierungsreihenfolgen:

\begin{sourcecode}[Anwendung]{csharp}
options.UseNpgsql(connectionString).UseAcmeTenancy().UseSnakeCaseNamingConvention();
\end{sourcecode}

\begin{sourcecode}[Erzeugtes Migrations-SQL (\allowbreak{}Auszug,\allowbreak{} tests/\allowbreak{}.\allowbreak{}.\allowbreak{}.\allowbreak{}/\allowbreak{}Naming\allowbreak{}Convention\allowbreak{}Tests.\allowbreak{}cs)]{sql}
ALTER TABLE orders ENABLE ROW LEVEL SECURITY;
CREATE POLICY tenant_isolation ON orders AS PERMISSIVE FOR ALL
    USING (tenant_id = nullif(current_setting('app.tenant_id', true), '')::uuid)
    WITH CHECK (tenant_id = nullif(current_setting('app.tenant_id', true), '')::uuid);
CREATE INDEX ix_orders_tenant_id ON orders (tenant_id);
\end{sourcecode}

\begin{tip}[Namen nie selbst zusammensetzen]
  Ein Plugin, das Namen selbst zusammensetzt (aus \code{Name}-Properties oder per
  String-Verkettung), bricht, sobald ein Naming-Plugin aktiv ist. Immer das Modell nach dem
  endgültigen Namen fragen, und zwar so spät wie möglich.
\end{tip}

\section{Attribut-Conventions}

Data Annotations wertet man mit einer Attribut-Convention aus. EF bietet dafür die
Basisklasse
\code{Type\allowbreak{}Attribute\allowbreak{}Convention\allowbreak{}Base<\allowbreak{}TAttribute>},
die das Attribut findet und nur noch die Reaktion verlangt. Der Wert wird mit
\code{from\allowbreak{}Data\allowbreak{}Annotation: true} gesetzt, sodass die Fluent API ihn
überstimmen kann:

\begin{sourcecode}[Conventions/\allowbreak{}No\allowbreak{}Tenant\allowbreak{}Isolation\allowbreak{}Attribute\allowbreak{}Convention.\allowbreak{}cs]{csharp}
public class NoTenantIsolationAttributeConvention
    : TypeAttributeConventionBase<NoTenantIsolationAttribute>
{
    public NoTenantIsolationAttributeConvention(
        ProviderConventionSetBuilderDependencies dependencies)
        : base(dependencies)
    {
    }

    protected override void ProcessEntityTypeAdded(
        IConventionEntityTypeBuilder entityTypeBuilder,
        NoTenantIsolationAttribute attribute,
        IConventionContext<IConventionEntityTypeBuilder> context)
        => entityTypeBuilder.HasRowLevelSecurity(enabled: false, fromDataAnnotation: true);
}
\end{sourcecode}

\begin{tip}
  Das Attribut selbst liegt im \emph{Kernpaket}, die Convention im Provider-Paket.
  Entitätsklassen können das Attribut verwenden, ohne von PostgreSQL abzuhängen. Das ist
  dieselbe Aufteilung wie bei EFs eigenen Attributen in
  \code{System.\allowbreak{}Component\allowbreak{}Model.\allowbreak{}Data\allowbreak{}Annotations}.
\end{tip}

% ------------------------------------------------------------
\chapter{Modellvalidierung}\label{ch:model-validation}

Ein Plugin sollte Fehlkonfiguration erkennen, bevor sie zu falschem SQL oder Datenlecks
führt. EF validiert Modelle über \code{IModel\allowbreak{}Validator}, den es aber nur einmal
gibt und der dem Provider gehört. Für Plugins ist eine
\code{IModel\allowbreak{}Finalized\allowbreak{}Convention} der bessere Ort: Sie läuft auf
dem fertigen Modell und kollidiert mit niemandem.

\begin{sourcecode}[Conventions/\allowbreak{}Tenancy\allowbreak{}Model\allowbreak{}Validation\allowbreak{}Convention.\allowbreak{}cs]{csharp}
public virtual IModel ProcessModelFinalized(IModel model)
{
    foreach (var entityType in model.GetEntityTypes())
    {
        if (!entityType.HasRlsConfiguration())
            continue;

        if (entityType.BaseType is not null)
            throw new InvalidOperationException(
                $"RLS an '{entityType.DisplayName()}' ist ungültig: nur am Root-Typ "
                + $"('{entityType.GetRootType().DisplayName()}') erlaubt.");

        if (StoreObjectIdentifier.Create(entityType, StoreObjectType.Table) is null)
            throw new InvalidOperationException("... nicht auf eine Tabelle gemappt ...");

        if (!entityType.IsRlsEnabled() && entityType.GetRlsPolicies().Count > 0)
            _logger.PoliciesWithoutRowLevelSecurity(_loggingDefinitions, entityType);
    }
    return model;
}
\end{sourcecode}

Die Validierung unterscheidet zwei Arten von Problemen. \textbf{Strukturelle Fehler}, die
garantiert zu falschem Verhalten führen, werfen eine Exception. \textbf{Fragwürdige
Konfiguration}, die gewollt sein kann, erzeugt eine Warnung über das Logging-System. Nutzer
können solche Warnungen gezielt in Fehler umwandeln
(Kapitel~\ref{ch:diagnostics-and-logging}).

% ------------------------------------------------------------
\chapter{Type Mappings und Value Converter}\label{ch:type-mappings-and-value-converters}

Ein Type Mapping verbindet einen CLR-Typ mit einem Datenbanktyp. Es bestimmt den Spaltentyp
in Migrationen, wie Parameter gesetzt werden, wie Literale im SQL aussehen, wie Werte
verglichen und kopiert werden und wie sie in JSON-Spalten gespeichert werden. Für eigene
Typen hat ein Plugin zwei Möglichkeiten:

\begin{datatable}{H{0.92}H{1.04}H{1.04}}
  \textbf{Weg} & \textbf{Vorteile} & \textbf{Nachteile} \\
  \tablehead
  Value Converter (Nutzer konfiguriert \code{Has\allowbreak{}Conversion}) & Kein Plugin
  nötig & Muss für jede Property oder per \code{Configure\allowbreak{}Conventions}
  konfiguriert werden \\
  \tablerule
  \code{IRelational\allowbreak{}Type\allowbreak{}Mapping\allowbreak{}Source\allowbreak{}Plugin}
  & Typ funktioniert überall ohne Konfiguration, eigene Literale, volle Kontrolle &
  Provider-spezifisch (Store-Typ) \\
\end{datatable}

\section{Der Typ und sein Konverter}

Der stark typisierte \code{Tenant\allowbreak{}Id} und sein Konverter liegen im Kernpaket.
Der Konverter hat eine gemeinsame Instanz; warum, zeigt der Abschnitt über Compiled Models.

\begin{sourcecode}[Tenant\allowbreak{}Id.\allowbreak{}cs,\allowbreak{} Tenant\allowbreak{}Id\allowbreak{}Converter.\allowbreak{}cs (\allowbreak{}gekürzt)]{csharp}
public readonly record struct TenantId(Guid Value)
{
    public static TenantId Empty => default;
    public bool IsEmpty => Value == Guid.Empty;              // wird nach SQL übersetzt
    public static TenantId New() => new(Guid.CreateVersion7());
}

public sealed class TenantIdConverter : ValueConverter<TenantId, Guid>
{
    public static TenantIdConverter Instance { get; } = new();

    public TenantIdConverter()
        : base(id => id.Value, value => new TenantId(value))
    {
    }
}
\end{sourcecode}

\section{Das Type Mapping}

\begin{sourcecode}[Storage/\allowbreak{}Tenant\allowbreak{}Id\allowbreak{}Type\allowbreak{}Mapping.\allowbreak{}cs]{csharp}
public class TenantIdTypeMapping : RelationalTypeMapping
{
    // Pflicht für Compiled Models: öffentliche statische Default-Property
    public static TenantIdTypeMapping Default { get; } = new();

    public TenantIdTypeMapping()
        : base(new RelationalTypeMappingParameters(
            new CoreTypeMappingParameters(
                typeof(TenantId),
                converter: TenantIdConverter.Instance,
                jsonValueReaderWriter: TenantIdJsonReaderWriter.Instance),
            storeType: "uuid",
            dbType: System.Data.DbType.Guid))
    {
    }

    protected TenantIdTypeMapping(RelationalTypeMappingParameters parameters)
        : base(parameters)
    {
    }

    // EF klont Mappings, z. B. für Facetten oder geänderte Comparer
    protected override RelationalTypeMapping Clone(
        RelationalTypeMappingParameters parameters)
        => new TenantIdTypeMapping(parameters);

    protected override string GenerateNonNullSqlLiteral(object value)
    {
        var guid = value switch
        {
            TenantId id => id.Value,
            Guid g => g,
            _ => throw new ArgumentException($"Unerwarteter Typ {value.GetType()}."),
        };
        return $"'{guid}'::uuid";
    }
}
\end{sourcecode}

Jedes eigene Mapping braucht den geschützten Konstruktor mit
\code{Relational\allowbreak{}Type\allowbreak{}Mapping\allowbreak{}Parameters} und eine
\code{Clone}-Überschreibung, die den eigenen Typ erzeugt. Fehlt sie, verliert ein geklontes
Mapping sein Verhalten, zum Beispiel die Literal-Erzeugung.

\section{Das Type-Mapping-Plugin}

\code{Relational\allowbreak{}Type\allowbreak{}Mapping\allowbreak{}Source} befragt alle
Plugins, \emph{bevor} es auf die Mappings des Providers zurückgreift. Ein Plugin gibt
\code{null} zurück, wenn es nicht zuständig ist:

\begin{sourcecode}[Storage/\allowbreak{}Tenancy\allowbreak{}Type\allowbreak{}Mapping\allowbreak{}Source\allowbreak{}Plugin.\allowbreak{}cs]{csharp}
public class TenancyTypeMappingSourcePlugin : IRelationalTypeMappingSourcePlugin
{
    public virtual RelationalTypeMapping? FindMapping(
        in RelationalTypeMappingInfo mappingInfo)
    {
        if (mappingInfo.ClrType != typeof(TenantId))
            return null;

        // Explizit anderer Spaltentyp (HasColumnType("text")): nicht zuständig
        var storeTypeBase = mappingInfo.StoreTypeNameBase;
        if (storeTypeBase is not null
            && !string.Equals(storeTypeBase, "uuid", StringComparison.OrdinalIgnoreCase))
            return null;

        return TenantIdTypeMapping.Default;
    }
}
\end{sourcecode}

\begin{note}
  Ohne Mapping würde EF \code{Tenant\allowbreak{}Id} beim Modellbau nicht als Skalartyp
  erkennen und mit einer Fehlermeldung abbrechen. Mit dem Plugin funktioniert der Typ als
  Property, als Schlüssel, in Abfragen, als Parameter und in JSON-Spalten, ohne eine Zeile
  Konfiguration.
\end{note}

\section{JSON-Unterstützung}

Für JSON-Spalten braucht ein Mapping einen
\code{Json\allowbreak{}Value\allowbreak{}Reader\allowbreak{}Writer}. Er hat drei Pflichten:
lesen, schreiben und beschreiben, wie er im generierten Code erzeugt wird. Das Muster folgt
EFs eigenem \code{Json\allowbreak{}Guid\allowbreak{}Reader\allowbreak{}Writer}:

\begin{sourcecode}[Tenant\allowbreak{}Id\allowbreak{}Json\allowbreak{}Reader\allowbreak{}Writer.\allowbreak{}cs]{csharp}
public sealed class TenantIdJsonReaderWriter : JsonValueReaderWriter<TenantId>
{
    private static readonly PropertyInfo InstanceProperty =
        typeof(TenantIdJsonReaderWriter).GetProperty(nameof(Instance))!;

    public static TenantIdJsonReaderWriter Instance { get; } = new();

    private TenantIdJsonReaderWriter()
    {
    }

    public override TenantId FromJsonTyped(
        ref Utf8JsonReaderManager manager, object? existingObject = null)
        => new(manager.CurrentReader.GetGuid());

    public override void ToJsonTyped(Utf8JsonWriter writer, TenantId value)
        => writer.WriteStringValue(value.Value);

    // Für Compiled Models und vorkompilierte Abfragen
    public override Expression ConstructorExpression
        => Expression.Property(null, InstanceProperty);
}
\end{sourcecode}

\begin{important}[Compiled Models brauchen Default und Instance]
  Beim Generieren für NativeAOT (\code{dotnet ef dbcontext optimize -{}-nativeaot}) erzeugt
  der Generator für Compiled Models
  (Kapitel~\ref{ch:compiled-models-trimming-and-nativeaot}) für jedes Type Mapping Code der
  Form
  \code{Tenant\allowbreak{}Id\allowbreak{}Type\allowbreak{}Mapping.\allowbreak{}Default.\allowbreak{}Clone(\allowbreak{}.\allowbreak{}.\allowbreak{}.\allowbreak{})}.
  Er verlangt eine \textbf{öffentliche, statische Property} namens \code{Default} auf einem
  \textbf{öffentlichen} Typ und bricht sonst ab. Konverter und Reader/Writer übernimmt er
  nur dann ohne zusätzlichen Code, wenn sie \emph{dieselbe Instanz} wie im Default-Mapping
  sind. Deshalb nutzt das Mapping die statischen Instanzen statt \code{new
  .\allowbreak{}.\allowbreak{}.\allowbreak{}(\allowbreak{})}. Ohne NativeAOT löst EF 10 Type
  Mappings stattdessen beim Start auf; ein Plugin sollte die Anforderungen trotzdem
  erfüllen, denn über NativeAOT entscheiden die Anwender, nicht das Plugin.
\end{important}
